\documentclass[12pt]{article}
\usepackage[utf8]{inputenc}
\usepackage[T1]{fontenc}
\usepackage{hyperref}
\title{Ray\'on Zoque Language Resources for K-12 Curriculum Development}
\author{Compiled by Claude Sonnet 5.5 (claude-sonnet-5-5) for Andr\'as Kornai}
\date{}
\begin{document}
\maketitle

\section{Introduction}
Ray\'on Zoque (ISO 639-3: zor; Glottocode rayo1235) is a Mixe-Zoquean language of Chiapas, Mexico, spoken in Ray\'on and Tapilula. One source estimates about 2,000 speakers. Glottolog lists it as shifting. Searches found one dictionary and a grammar of a related variety, and no school materials.

\section{Primary Educational Resources}
No primers, textbooks or school programmes were found. This section is empty.

\section{Linguistic and Reference Materials}
\textit{Diccionario zoque de Ray\'on} (Harrison, Harrison, L\'opez Ju\'arez and Orodo\~nes, 1984, Summer Institute of Linguistics, Mexico) is cited in the Glottolog entry for the language. \url{https://glottolog.org/resource/languoid/id/rayo1235}

A 1982 \textit{Diccionario del dialecto Zoque} by Porfirio Cundap\'i Estrada is catalogued with a note that the variety is probably Ray\'on Zoque. This is only a tentative identification. \url{https://glottolog.org/resource/reference/id/318605}

Faarlund's \textit{A Grammar of Chiapas Zoque} (Oxford University Press, 2012) describes the Ocotepec and Tapalapa variety. Ray\'on is in the same northeast dialect area, but this is not a description of Ray\'on itself. \url{https://en.wikipedia.org/wiki/Chiapas_Zoque}

\section{Religious and Community Resources}
None found.

\section{Digital and Audio Resources}
A numerals page records Ray\'on Zoque numeral data. \url{https://lingweb.eva.mpg.de/channumerals/Zoque-Rayon.htm}

\section{Curriculum Development Context}
The language is thought to be used by adults only and is not known to be taught in schools. The 1984 dictionary is the main written source. A curriculum would have to build primers and reading material from it and from the related Chiapas Zoque description. Both are secondary to work with speakers.

\end{document}